\documentclass[12pt]{article}
\usepackage[utf8]{inputenc}
\usepackage{geometry}
\geometry{margin=1in}
\usepackage{setspace}
\usepackage{lipsum}

\title{Témoignage de Philippe Thomas Savard}
\author{}
\date{}

\begin{document}

\maketitle

\onehalfspacing

Témoignage : Philippe Thomas Savard  
Né le 27 novembre 1984 à Québec  
Adresse actuelle du domicile : 4-5354 rue du Menuet G6X 2Y6  
Numéro de téléphone : 418-487-8798 et 418-832-4531  
Emplois et occupation : Installateur de piscine hors terre chez Concept Piscine Désigné. Développeur de validation mathématique sur Isabelle code HOL / conjecture de la fonction zêta de Riemann.  
\textbf{Inscrit et admis en dessin industriel au CFP Neufchâtel}

Témoignage :  
Les policiers sont venus à la maison (4-5354 rue du Menuet à Lévis) où je demeure. Ils ont cogné et j’ai répondu à la demande des policiers en leur ouvrant la porte d’entrée de mon domicile. Ils mentionnent qu’ils ont un mandat et qu’ils viennent procéder à mon arrestation et à la perquisition de mon domicile. Ils entrent, les menottes me sont mises aux poignets mais retirées après quelques minutes, 4 ou 5 minutes plus tard, et une fois m’avoir informé que j’étais mis en état d’arrestation et que mes droits m’ont été mentionnés. Ils m’offrent de subir un interrogatoire à deux endroits différents selon mes désirs : poste de police rue du Sault ou dans mon appartement directement. J’ai décidé de suivre les policiers jusqu’au poste de police rue du Sault.

Au poste de police, durant que la perquisition était en cours, deux enquêteurs sont venus me rejoindre dans une pièce réservée à cet effet et m’ont lu de nouveau mes droits. Ils m’ont fait signer un document pour s’assurer que je comprenais que j’avais le droit de garder le silence et que ma signature ne pouvait servir à autre chose. J’ai signé le document attestant que je comprenais que j’avais le droit de garder le silence.

Ils m’ont ensuite proposé de contacter un avocat et que si je n’en connaissais pas, ils détenaient la liste des avocats de l’aide juridique. J’ai pu parler à une avocate (le nom m’échappe) qui m’informe que j’ai le droit de garder le silence et que je ne suis pas tenu de répondre aux questions des policiers. Elle me le conseille fortement, que son expérience d’avocate conseille de ne répondre à aucune question des policiers. Que si je le fais, tout ce que je mentionne aux policiers peut être retenu contre moi au tribunal si des accusations sont portées. Elle m’explique que les policiers sont des professionnels et qu’ils peuvent utiliser des techniques d’interrogation en lien avec leur nature professionnelle, qu’ils ne sont pas mes amis. Je tiens à affirmer que les policiers sont les bienvenus chez moi et que le travail qu’ils effectuent m’apparaît nécessaire à notre collectivité et au bien-être de chacun.

Elle me spécifie clairement que dans les méthodes d’interrogation, les policiers sont en droit d’affirmer de manière paracohérente des informations vrai/faux et faux/vrai pour m’interroger. J’ai donc précisé aux policiers durant l’interrogatoire, à chacune des questions qui m’étaient dirigées, en mentionnant à chaque occasion : « Je vais garder le silence, merci ».

Je voudrais spécifier par l’affirmation suivante : « qu’il est évident qu’il n’y a pas, depuis que je demeure à l’adresse de mon domicile, eu que de très rares fois des gens âgés d’âge mineur à mon domicile ; qu’il s’agit de mes neveux et nièces qui sont accompagnés toujours de leurs parents. Je suis un homme vivant seul, sans enfant et sans conjointe depuis toujours, et qu’il n’y a à toute fin pratique jamais de mineurs chez moi ».

Par suite de l’interrogatoire, les policiers m’ont conduit au poste de police de Lévis centre-ville sur le boulevard Guillaume-Couture. De là, j’ai été libéré et on m’a dit que quelqu’un était dans le stationnement du poste de police pour assurer mon transport jusqu’à mon domicile. Je suis retourné à mon domicile par mes propres moyens. Ces événements ont eu lieu en février 2025. À la mi-juin 2025, par la poste, j’ai reçu l’ordre par sommation de me présenter au tribunal.

Déclaration de l’inspecteur Gauthier lors de mon arrestation :  
– En s’adressant à moi : « Nous les policiers, on invente les accusations qui sont portées contre toi ».  
– En parlant de mon domicile et à mon domicile, en s’adressant à moi : « Ce n’est de toute façon pas chez toi ici ».

Propos tenus par l’homme assurant la maintenance de l’immeuble où je demeure, Benoît Charette :  
En parlant du locataire de mon adresse juste avant le début du nouveau bail le 1er juillet 2025 : « Le propriétaire de l’immeuble te fait dire que si tu veux retirer le nom de ton colocataire qui reste avec toi, t’as juste à le mentionner au propriétaire de l’immeuble ».

Affirmation de ma part sur l’affirmation de l’homme assurant la maintenance de l’immeuble :  
« J’ai un titre de propriété lié à un droit personnel lié à un bail de location à l’adresse de mon domicile et le bail est à mon nom seul ».  
Affirmation de ma part sur l’affirmation de l’homme assurant la maintenance de l’immeuble :  
« Je reçois régulièrement du courrier à mon adresse portant le nom d’un Philippe Savard et qui ne m’est pas adressé. De la Croix-Rouge, organisme avec lequel je n’ai aucun lien ni réseau établi. Des sociétés pour les malades qui font des levées de fonds pour financer leurs causes, avec lesquelles je n’ai aucun lien ni réseau ».

Situation à mon domicile rapportée et affirmée :  
« Certaines de mes fréquentations comme Robert Mineault et Cathie Talbot viennent à la maison. Robert (Bob) Mineault est sans domicile fixe et vient régulièrement se réfugier chez moi. Selon les informations, il est connu des policiers. Cathie Talbot, avec qui je crois qu’il entretient une relation intime, est connue selon les informations que je détiens en lien avec la consommation de drogue. Je suis un homme solitaire, je suis toujours à la maison seul, je ne sors de la maison qu’à de très rares occasions. Quand Robert Mineault vient héberger à la maison, je suis dans ma chambre à écouter la télé et lui reste dans la pièce du salon, vu ma nature solitaire. Je sais, pour l’avoir vu, que Cathie vient lui rendre visite pendant que je suis dans la chambre seul, la porte close.  
J’ai une prescription à mon dossier au Brunet Mont-Marie pour soigner des problèmes érectiles ».

--- Fin de mon témoignage.

\vspace{1cm}

\noindent\rule{8cm}{0.4pt}\\
Philippe Thomas Savard

\end{document}
